\documentclass{beamer}

\usepackage[utf8]{inputenc}
\usepackage[spanish]{babel}
\usepackage{tikz}
\usepackage{booktabs}
\usepackage{pgfplots}
% \usepackage{titling} % adds \theauthor \thetitle etc
\usepgfplotslibrary{dateplot}

\newcommand{\titulo}{Git: software de control de versiones}
\newcommand{\autor}{Daniel Lubián Arenillas}
\newcommand{\fecha}{\today}
\newcommand{\git}{\texttt{git}}

\title{\titulo}
\author{\autor}
\institute[IDR/UPM]{IDR/UPM}
\date{\fecha}
% \logo{\includegraphics[height=1.5cm]{fig/tres.pdf}}



% \usetheme{Rochester}
\usetheme{metropolis}
\metroset{block=fill, numbering=fraction, sectionpage=progressbar,subsectionpage=progressbar}%,titleformat=smallcaps}
% \usecolortheme{structure}

% \setbeamertemplate{frame footer}{\autor\,--\,\titulo}
% \definecolor{azuletsiae}{RGB}{59,80,141}
% \definecolor{azulclaroetsiae}{RGB}{197,208,228}
% \definecolor{black1}{RGB}{26,28,34}
% \setbeamercolor{normal text}{fg=black1}
% \setbeamercolor{frametitle}{bg=azuletsiae}
% % \setbeamercolor{frametitle}{bg=black1}
% \setbeamercolor{progress bar}{fg=black1}
% \setbeamercolor{progress bar}{fg=azuletsiae,bg=white}
% \setbeamercolor{progress bar}{fg=azuletsiae,bg=azulclaroetsiae}
% \titlegraphic{
	% \includegraphics[width=2cm]{fig/idr.pdf}
	% \hfill
	% \includegraphics[width=2cm]{fig/tres.pdf}
% }

% \usepackage{helvet}
\usepackage[sfdefault]{FiraSans} %% option 'sfdefault' activates Fira Sans as the default text font
\usepackage{FiraMono}
% \usepackage[sfdefault,light,condensed]{roboto}
\usepackage[T1]{fontenc}




% \AtBeginSection[] % add toc at the beginning of a section, and highlight the next one
% {\begin{frame}
% 		\frametitle{Table of Contents}
% 		\tableofcontents[currentsection]
% \end{frame}}

% \addtobeamertemplate{frametitle}{}{%
% 	\begin{tikzpicture}[remember picture,overlay]
% 	\node[anchor=north east,yshift=1pt] at (current page.north east) {\includegraphics[height=1.4cm]{fig/tres}};
% 	\end{tikzpicture}}

\begin{document}
	
\maketitle

% \begin{frame}\frametitle{Índice}
% 	\tableofcontents
% \end{frame}

\begin{frame}\frametitle{¿Por qué?}
	\centering
	\includegraphics[height=8cm]{fig/phd.jpeg}
\end{frame}

\begin{frame}\frametitle{¿Qué es \git?}
	\begin{columns}
		\begin{column}{0.5\textwidth}
			\begin{block}{Wikipedia}
				\git\ es un software de control de versiones diseñado por Linus Torvalds, pensando en la eficiencia y la confiabilidad del mantenimiento de versiones de aplicaciones cuando éstas tienen un gran número de archivos de código fuente
			\end{block}
		\end{column}
		\begin{column}{0.5\textwidth}
			\centering
			\includegraphics[width=0.9\textwidth]{fig/git}
		\end{column}
	\end{columns}
\end{frame}

\begin{frame}\frametitle{¿Para qué puede sernos útil?}
	Es especialmente potente con archivos de texto
	\begin{itemize}
		\item Trabajos con Matlab, Fortran, Python, C, GMAT ... cualquier código fuente
		\item Documentos hechos con \LaTeX
		\item Cualquier otro archivo de texto (p.ej: datos de resultados): \texttt{.txt}, \texttt{.csv}, \texttt{.dat} ...
	\end{itemize}
	Con archivos distintos pierdes funcionalidades, por lo que pierde algo de sentido. Requiere ser mucho más organizado
	\begin{itemize}
		\item Catia, Excel, Word,...
	\end{itemize}
\end{frame}

\begin{frame}\frametitle{Comandos más importantes}
	\begin{itemize}
		\item \textbf{\texttt{init}		}: inicializa el VCS
		\item \textbf{\texttt{status}	}: cómo están las cosas $\rightarrow$ \underline{estado}
		\item \textbf{\texttt{add}		}: \underline{añade} archivos al commit
		\item \textbf{\texttt{commit}	}: \underline{comete} (hace efectivos) esos cambios
		\item \textbf{\texttt{diff}		}: muestra las \underline{diferencias}
		\item \textbf{\texttt{blame}	}: muestra el ``\underline{culpable}'' de cada línea
		\item \textbf{\texttt{branch}	}: crea/maneja \underline{ramas} del código
		\item \textbf{\texttt{merge}	}: \underline{mezcla} commits, branches, etc
		\item \textbf{\texttt{checkout}	}: \underline{cambia a} otro commit u otra rama
	\end{itemize}
\end{frame}

\begin{frame}\frametitle{Comandos más importantes (repos remotos)}
	\begin{itemize}
		\item \textbf{\texttt{clone}	}: crea una \underline{copia} de un repo remoto en la que estás
		\item \textbf{\texttt{push}		}: ``empuja'' (\underline{sube}) los commits del repo local a un repo remoto
		\item \textbf{\texttt{pull}		}: actualiza (\underline{baja}) los commits del repo remoto en el local
	\end{itemize}
\end{frame}

\begin{frame}
	\begin{block}{\texttt{git status}}
		
	\end{block}
\end{frame}


\begin{frame}\frametitle{Repositorios remotos (I)}

\end{frame}

\begin{frame}\frametitle{Repositorios remotos (II)}
	\begin{columns}
		\begin{column}{0.5\textwidth}
			\begin{center}
				\includegraphics[height = 1.3 cm]{fig/github/GitHub-Mark-120px-plus.png}
				\includegraphics[height = 1.3 cm]{fig/github/GitHub_Logo.png}
			\end{center}
			\begin{itemize}
				\item \url{https://github.com/}
				\item Código cerrado
				\item Sin pagar \textbf{no} hay repositorios privados
				\item Mucho más conocido
			\end{itemize}
		\end{column}
		\begin{column}{0.5\textwidth}
			\begin{center}
				\includegraphics[height = 1.3 cm]{fig/gitlab_wm_no_bg.pdf}
			\end{center}
			\begin{itemize}
				\item \url{https://gitlab.com/}
				\item Código abierto
				\item Sin pagar \textbf{sí} hay repositorios privados
				\item ¿?
			\end{itemize}
		\end{column}
	\end{columns}

\vspace{0.5cm}

\textbf{Otros servicios:} Bitbucket (\url{https://bitbucket.org/})

\end{frame}

\end{document}
